\Needspace{10\baselineskip}
\begin{center}\emph{The residues $+5$ and $-5$.}\end{center}

\Article{121}
Inspection suggests that $+5$ is a nonresidue of every odd number of the form $5n+2$ or $5n+3$, \emph{that is,} of every odd number that is a nonresidue of $5$. We prove that this rule has no exceptions as follows. Let the smallest exception, if there is one, be $=t$. It is therefore a nonresidue of $5$, while $5$ is a residue of $t$. Let $aa=5+tu$, with $a$ even and smaller than $t$. Then $u$ is odd and smaller than $t$, and $+5$ is a residue of $u$. If $a$ is not divisible by $5$, neither is $u$. But $tu$ is clearly a residue of $5$, so, since $t$ is a nonresidue of $5$, $u$ is a nonresidue too. \emph{That is,} there is an odd nonresidue of $5$, of which $+5$ is a residue, smaller than $t$, contrary to the hypothesis. If $a$ is divisible by $5$, \SourceLeaf{96}put $a=5b$ and $u=5v$; then $tv\equiv -1\equiv 4\,(\operatorname{mod.}5)$, \emph{that is,} $tv$ is a residue of 5. The rest of the proof proceeds as in the first case.

\Article{122}
Thus, for all primes that are nonresidues of 5 and have the form $4n+1$, \emph{that is,} all primes of the form $20n+13$ or $20n+17$, both $+5$ and $-5$ are nonresidues. For all primes of the form $20n+3$ or $20n+7$, $+5$ is a nonresidue and $-5$ a residue.

In just the same way we can prove that $-5$ is a nonresidue of all primes of the forms $20n+11$, $20n+13$, $20n+17$, $20n+19$. It follows that $+5$ is a residue of all primes of the form $20n+11$ or $20n+19$, and a nonresidue of all those of the form $20n+13$ or $20n+17$. Every prime except 2 and 5 (of which $\pm 5$ is a residue) has one of the forms $20n+1,3,7,9,11,13,17,19$, so we can now decide every case except those of the form $20n+1$ or $20n+9$.

\Article{123}
Inspection readily suggests that $+5$ and $-5$ are residues of all primes of the form $20n+1$ or $20n+9$. If this holds in general, we have the elegant law that \emph{$+5$ is a residue of every prime that is a residue of 5} (these have one of the forms $5n+1$ or $5n+4$, or one of $20n+1,9,11,19$; the third and fourth of the latter have already been dealt with), \emph{but a nonresidue of every odd number that is a nonresidue of 5}, as we proved above. This theorem is enough to decide whether $+5$ (and therefore $-5$, considered as the product of $+5$ and $-1$) is a residue or a nonresidue of any given number. Finally, note the analogy between this theorem and the one concerning the residue $-3$ in Article 120.

But proving this observed pattern is not so easy. For a prime of the form $20n+1$, or more generally of the form $5n+1$, we can proceed much as in Articles 114 and 119. For the modulus $5n+1$, let $a$ be any number belonging to the exponent 5; such numbers exist by the previous section. \SourceLeaf{97}Then $a^5\equiv1$, or $(a-1)(a^4+a^3+a^2+a+1)\equiv0\ (\mathrm{mod.}\,5n+1)$. Since we cannot have $a\equiv1$, and therefore cannot have $a-1\equiv0$, we must have $a^4+a^3+a^2+a+1$ $\equiv0$. Thus $4(a^4+a^3+a^2+a+1)=(2aa+a+2)^2-5a^2$ is also $\equiv0$, \textit{that is,} $5a^2$ is a residue of $5n+1$, and so is $5$, since $a^2$ is a residue not divisible by $5n+1$ (for $a$ is not divisible by $5n+1$ because $a^5\equiv1)$. \mbox{\textit{Q. E. D.}}

The case of a prime of the form $5n+4$, however, requires subtler methods. The propositions used to resolve it will be treated more generally later, so we touch on them only briefly here.

\Needspace{8\baselineskip}
I. If $p$ is a prime and $b$ a given quadratic nonresidue of $p$, the value of the expression

\[
(A)\ldots\quad\frac{(x+\surd b)^{p+1}-(x-\surd b)^{p+1}}{\surd b}
\]

(from whose expansion the irrationality clearly disappears) is always divisible by $p$, whatever number is substituted for $x$. Inspection of the coefficients in the expansion of $A$ shows that all terms from the second to the next-to-last, inclusive, are divisible by $p$, so $A\equiv2(p+1)(x^p+xb^{\frac{p-1}{2}})$ $(\mathrm{mod.}\,p)$. Since $b$ is a nonresidue of $p$, we have $b^{\frac{p-1}{2}}\equiv-1\ (\mathrm{mod.}\,p)$ (Article $106$), while $x^p$ is always $\equiv x$ (previous section). So $A\equiv0$. \mbox{\textit{Q. E. D.}}

\Needspace{8\baselineskip}
II. In the congruence $A\equiv0\ (\mathrm{mod.}\,p)$, the indeterminate $x$ has degree $p$, and all the numbers $0,1,2....p-1$ are roots. Let $e$ be a divisor of $p+1$. Then the expression

\[
\frac{(x+\surd b)^e-(x-\surd b)^e}{\surd b}
\]

(which we denote by $B$) is free of irrationality when expanded, and the indeterminate $x$ has degree $e-1$ in it. Elementary analysis tells us that $A$ is divisible by $B$ identically. I now say that there are $e-1$ values of $x$ which, substituted into $B$, make $B$ divisible by $p$. For put $A\equiv BC$. Then $x$ in $C$ has degree $p-e+1$, so the congruence $C\equiv0\ (\mathrm{mod.}\,p)$ has at most $p-e+1$ roots. Thus all the remaining numbers among $0,1,2,3....p-1$, whose number is $=e-1$, are roots of the congruence $B\equiv0$.

\Needspace{8\baselineskip}
III. Now let $p$ have the form $5n+4$, let $e=5$, let $b$ be a nonresidue of $p$, and determine $a$ so that

\SourceLeaf{98}\[
\frac{(a+\surd b)^5-(a-\surd b)^5}{\surd b}
\]

is divisible by $p$. This expression becomes

\[
=10a^4+20aab+2bb=2((b+5aa)^2-20a^4)
\]

So $(b+5aa)^2-20a^4$ too is divisible by $p$, \textit{that is,} $20a^4$ is a residue of $p$. Since $4a^4$ is a residue not divisible by $p$ (we readily see that $a$ cannot be divisible by $p$), $5$ is also a residue of $p$. \mbox{\textit{Q. E. D.}}

Thus the theorem stated at the beginning of this article holds in general. ---

We also note that the proofs for both cases are due to the distinguished Lagrange, \textit{M\'em. de l'Ac. de Berlin}, $1775$, \textit{p.} $352$ ff.

\Needspace{10\baselineskip}
\begin{center}\textit{On $\pm7$.}\end{center}

\Article{124}
By a similar method we prove that

$-7$ \textit{is a nonresidue of every number that is a nonresidue of $7$.}

Inspection suggests, on the other hand, that

$-7$ \textit{is a residue of every prime that is a residue of $7$.}

But no one has yet proved this rigorously. For those residues of $7$ that have the form $4n-1$, the proof is easy: a method already familiar from what precedes shows that $+7$ is always a nonresidue of such primes, and so $-7$ a residue. But this gains us little, since the remaining cases cannot be treated by this method. We can handle one more case much as in Articles $119$, $123$. If $p$ is a prime of the form $7n+1$, and $a$, for the modulus $p$, belongs to the exponent $7$, then

\[
\frac{4(a^7-1)}{a-1}=(2a^3+a^2-a-2)^2+7(a^2+a)^2
\]

is clearly divisible by $p$, so $-7(a^2+a)^2$ is a residue of $p$. But $(a^2+a)^2$, being a square, is a residue of $p$, and is not divisible by $p$. For since $a$ is assumed to belong to the exponent $7$, it can be neither $\equiv0$ nor $\equiv-1\ (\mathrm{mod.}\,p)$; \textit{that is,} neither $a$ nor $a+1$ is divisible by $p$, and neither is the square $(a+1)^2a^2$. Thus $7$ too is clearly a residue of $p$. \SourceLeaf{99}\mbox{\textit{Q. E. D.}} --- But primes of the form $7n+2$ or $7n+4$ resist all the methods given so far. This proof too was first found by the distinguished Lagrange at the cited place. --- In Section VII below we will show in general that the expression $\frac{4(x^p-1)}{x-1}$ can always be reduced to the form $X^2\mp pY^2$ (taking the upper sign when $p$ is a prime of the form $4n+1$, and the lower when it has the form $4n+3$), where $X,Y$ denote rational functions of $x$, free of fractions. Lagrange did not carry out this decomposition beyond the case $p=7$; see the cited place, p. $352$.

\Needspace{10\baselineskip}
\begin{center}\textit{Preparation for the general investigation.}\end{center}

\Article{125}
The preceding methods do not suffice to establish general proofs, so it is time to present another that is free of this defect. We begin with a theorem whose proof eluded our efforts for quite some time, although at first sight it seems so obvious that some have not even seen the need for a proof. It is this: \textit{Every number, except squares taken positively, is a nonresidue of some primes.} We will use this theorem only as an aid in proving others, so we treat here only the cases needed for that purpose. The remaining cases will follow of their own accord later. We will therefore show that \textit{every prime of the form $4n+1$, taken either positively or negatively}\HistoricalNote{*)}{The need to exclude $+1$ is clear.}, \textit{is a nonresidue of some primes}, and, if $>5$, of primes smaller than itself.

\textit{First}, suppose the prime $p$, of the form $4n+1$ ($>17$; but $-13\,N3$, $-17\,N5$), is to be taken \textit{negatively}. Let $2a$ be the even number immediately greater than $\surd p$. We readily see that $4aa$ is always $<2p$, or $4aa-p<p$. Now $4aa-p$ has the form $4n+3$, and $+p$ is a quadratic residue of $4aa-p$ (since $p\equiv4aa\ (\mathrm{mod.}\,4aa-p)$). If $4aa-p$ is prime, $-p$ is therefore a nonresidue of it. If not, $4aa-p$ must have a factor of the form $4n+3$; since $+p$ must be a residue of this factor too, $-p$ is a nonresidue of it. \mbox{\textit{Q. E. D.}}

For primes taken \textit{positively}, we distinguish two cases. \textit{First}, let $p$ be a prime of the form $8n+5$. Let $a$ be any positive number $<\surd\!\frac12 p$. Then $8n+5-2aa$ is a positive number of the form $8n+5$ or $8n+3$, according as $a$ \SourceLeaf{100}is even or odd. It must therefore be divisible by a prime of the form $8n+3$ or $8n+5$, since a product of any number of factors of the forms $8n+1$ and $8n+7$ cannot have either the form $8n+3$ or $8n+5$. Let this prime be $q$; then $8n+5\equiv2a^2\ (\mathrm{mod.}\,q)$. But $2$ is a nonresidue of $q$ (Article $112$), and therefore so are $2a^2$\HistoricalNote{*)}{Article $98$. For $a^2$ is clearly a residue of $q$ not divisible by $q$; otherwise the prime $p$ too would be divisible by $q$. \mbox{\textit{Q. E. A.}}} and $8n+5$. \mbox{\textit{Q. E. D.}}

\Article{126}
But that every prime of the form $8n+1$, taken positively, is a nonresidue of some smaller prime cannot be proved by such obvious devices. This truth is of the greatest importance, so we cannot pass over a rigorous proof, even though it is somewhat long. We first give the following

\Needspace{10\baselineskip}
\textsc{Lemma.} \textit{Suppose two sequences of numbers}

\[
A,B,C\ \textit{etc.}\ ....\ (\mathrm{I}),\quad A^{\prime},B^{\prime},C^{\prime}\ \textit{etc.}\ ....\ (\mathrm{II})
\]

(whether they have the same number of terms does not matter) \textit{are such that, for any prime or prime power $p$ dividing a term of the second sequence, or several terms, at least as many terms of the first sequence are divisible by $p$ as of the second. Then I say that the product of all the numbers} (I) \textit{is divisible by the product of all the numbers} (II).

\textit{Example.} Let (I) consist of $12,18,45$ and (II) of $3,4,5,6,9$. Then the numbers of terms divisible by $2,4,3,9,5$ are respectively $2,1,3,2,1$ in (I) and $2,1,3,1,1$ in (II). The product of all terms of (I), $=9720$, is divisible by the product of all terms of (II), $3240$.

\textit{Proof.} Let the product of all terms of (I) be $=Q$, and that of all terms of sequence (II) $=Q'$. Any prime divisor of $Q'$ clearly also divides $Q$. We now show that every prime factor of $Q'$ occurs in $Q$ with an exponent at least as large as in $Q'$. Let $p$ be such a divisor, and suppose that sequence (I) has $a$ terms divisible by $p$, $b$ terms divisible by $p^2$, $c$ terms divisible by $p^3$, etc. Let $a',b',c'$, etc. have analogous meanings for sequence (II). We readily see that $p$ occurs in $Q$ with exponent \SourceLeaf{101}$a+b+c+$, etc., and in $Q'$ with exponent $a'+b'+c'+$, etc. But $a'$ is no greater than $a$, $b'$ no greater than $b$, etc., by hypothesis; so $a'+b'+c'$, etc. is certainly not $>a+b+c$, etc. --- Thus no prime can have a larger exponent in $Q'$ than in $Q$, so $Q$ is divisible by $Q'$ (Article $17$). \mbox{\textit{Q. E. D.}}

\Needspace{10\baselineskip}
\Article{127}
\textsc{Lemma.} \textit{The progression $1,2,3,4,....n$ cannot have more terms divisible by any number $h$ than the progression $a,a+1,a+2.....a+n-1$, which has the same number of terms.}

If $n$ is a multiple of $h$, both progressions clearly have $\frac nh$ terms divisible by $h$. Otherwise, put $n=eh+f$, with $f$ $<h$. The first sequence then has $e$ terms divisible by $h$, and the second either just as many or $e+1$.

As a corollary we obtain the proposition known from the theory of figurate numbers, but, unless we are mistaken, never yet proved directly, that

\[
\frac{a.a+1.a+2....a+n-1}{1\ .\ 2\ .\ 3\ ....\ n}
\]

is always an integer.

Finally, the lemma could have been stated more generally as follows:

In the progression $a,a+1,a+2,....a+n-1$, there are at least as many terms congruent modulo $h$ to any given number $r$ as there are in $1,2,3....n$ divisible by $h$.

\Needspace{12\baselineskip}
\Article{128}
\textsc{Theorem.} \textit{Let $a$ be any number of the form $8n+1$, let $p$ be any number coprime to $a$ of which $+a$ is a residue, and finally let $m$ be arbitrary. Then I say that the progression}

\[
\begin{gathered}a,\ \tfrac12(a-1),\ 2(a-4),\ \tfrac12(a-9),\ 2(a-16),\\ ....\ 2(a-m^2),\quad\textit{or}\quad\tfrac12(a-m^2)\end{gathered}
\]

\textit{according as $m$ is even or odd, has at least as many terms divisible by $p$ as the progression}

\[
1,\ 2,\ 3,\ ....\ 2m+1
\]

We denote the first progression by (I) and the second by (II).

\SourceLeaf{102}\textit{Proof.} I. When $p=2$, all terms of (I) except the first, \textit{that is,} $m$ terms, are divisible; there are just as many in (II).

II. Let $p$ be odd, or twice or four times an odd number, and let $a\equiv rr\ (\mathrm{mod.}\,p)$. Then the progression $-m,-(m-1),-(m-2),.....+m$ (which has as many terms as $(\mathrm{II})$ and will be denoted by (III)) has at least as many terms congruent modulo $p$ to $r$ as sequence (II) has terms divisible by $p$ (preceding article). No two of them can differ only in sign and not in magnitude\HistoricalNote{*)}{For if $r\equiv-f\equiv+f\,(\operatorname{mod.}p)$, $2f$ would be divisible by $p$, and therefore so would $2a$ (because $ff\equiv a\,(\operatorname{mod.}p)$). This is possible only if $p=2$, since $a$ is coprime to $p$ by hypothesis. But we have already treated this case separately.}. Finally, each has a corresponding term of (I) divisible by $p$. For if $\pm b$ is a term of (III) congruent to $r$ modulo $p$, then $a-bb$ is divisible by $p$. If $b$ is even, the term $2(a-bb)$ of (I) is divisible by $p$. If $b$ is odd, the term $\frac12(a-bb)$ is divisible by $p$: for $\frac{a-bb}{p}$ is clearly an \textit{even} integer, since $a-bb$ is divisible by $8$, while the highest power of two dividing $p$ is at most $4$ (by hypothesis $a$ has the form $8n+1$, and $bb$, as the square of an odd number, has the same form, so their difference has the form $8n$). Thus sequence (I) has as many terms divisible by $p$ as (III) has terms congruent to $r$ modulo $p$, \textit{that is,} at least as many as (II) has terms divisible by $p$. \mbox{\textit{Q. E. D.}}

III. Let $p$ have the form $8n$, and let $a\equiv rr\ (\mathrm{mod.}\,2p)$. We readily see that $a$, being a residue of $p$ by hypothesis, is also a residue of $2p$. Then sequence (III) has at least as many terms congruent to $r$ modulo $p$ as (II) has terms divisible by $p$, and all of them differ in magnitude. But each has a corresponding term of (I) divisible by $p$. For if $+b$ or $-b\equiv r\ (\mathrm{mod.}\,p)$, then $bb\equiv rr\ (\mathrm{mod.}\,2p)$\HistoricalNote{\(\dagger\))}{For $bb-rr=(b-r)(b+r)$ is a product of two factors, one divisible by $p$ by hypothesis, the other by $2$ (since both $b$ and $r$ are odd); so $bb-rr$ is divisible by $2p$.}, so the term $\frac12(a-bb)$ is divisible by $p$. Thus (I) has at least as many terms divisible by $p$ as (II). \mbox{\textit{Q. E. D.}}

\Needspace{10\baselineskip}
\Article{129}
\textsc{Theorem.} \textit{If $a$ is a prime of the form $8n+1$, then below $2\surd a+1$ there must be a prime of which $a$ is a nonresidue.}

\SourceLeaf{103}\textit{Proof.} Suppose, if possible, that $a$ is a residue of all primes smaller than $2\surd a+1$. Then $a$ is also a residue of all composite numbers smaller than $2\surd a+1$ (compare the rules for deciding whether a number is a residue of a composite number, Article $105$). Let the integer immediately below $\surd a$ be $=m$. Then the sequence

\[
\begin{gathered}(\mathrm I)......a,\tfrac12(a-1),2(a-4),\tfrac12(a-9)....2(a-mm)\\ \text{or}\quad\tfrac12(a-mm)\end{gathered}
\]

has at least as many terms divisible by any number smaller than $2\surd a+1$ as the sequence

\[
(\mathrm{II}).....1,2,3,4....2m+1\quad\text{(preceding article)}
\]

So the product of all terms of (I) is divisible by the product of all terms of (II) (Article $126$). But the former is either \mbox{$=a(a-1)(a-4)....(a-mm)$} or half this product, according as $m$ is even or odd. Thus the product $a(a-1)(a-4)....(a-mm)$ is certainly divisible by the product of all terms of $(\mathrm{II})$; and since all these terms are coprime to $a$, this remains true when the factor $a$ is omitted. But the product of all terms of $(\mathrm{II})$ can also be written

\[
(m+1).((m+1)^2-1).((m+1)^2-4)....((m+1)^2-m^2)
\]

So

\[
\frac1{m+1}.\frac{a-1}{(m+1)^2-1}.\frac{a-4}{(m+1)^2-4}....\frac{a-m^2}{(m+1)^2-m^2}
\]

is an integer, even though it is a product of fractions smaller than one. For $\surd a$ must be irrational, so $m+1>\surd a$, and therefore $(m+1)^2>a.$ Thus our assumption is impossible. \mbox{\textit{Q. E. D.}}

Since $a$ is certainly $>9$, we have $2\surd a+1<a$, so there is a prime $<a$ of which $a$ is a nonresidue.
